\chapter{Instruments}
\label{app:instruments}

\section{The five describer questions}
The describer is asked, in this order, and its answers are stored per record:
\begin{enumerate}\setlength{\itemsep}{1pt}
\item \emph{What is the facial expression?} (\texttt{expression})
\item \emph{What is the person doing?} (\texttt{action})
\item \emph{What is the body posture?} (\texttt{posture})
\item \emph{How many people are there?} (\texttt{people})
\item \emph{What is in the picture?} (\texttt{scene})
\end{enumerate}
The word ``emotion'' appears in none of them.

\section{The description template}
The five answers are composed into one sentence:
\begin{quote}\ttfamily
A drawing of \{people\} showing \{scene\}. The person is \{action\}, \{posture\}. The facial expression is \{expression\}.
\end{quote}
The template is applied to every record, including those where \texttt{people} is ``0'' or ``none''; the ablation variants remove one clause at a time or fix \texttt{expression} to ``unclear''.

\section{The image reader's hypothesis}
Zero-shot classification by entailment with the hypothesis template
\begin{quote}\ttfamily
This scene conveys \{family\}.
\end{quote}
scored for each of the seven families; the three highest entailment scores are kept.

\section{GoEmotions labels to families}
\begin{description}\setlength{\itemsep}{1pt}
\item[joy] admiration, amusement, approval, caring, desire, excitement, gratitude, joy, love, optimism, pride, relief
\item[sadness] disappointment, embarrassment, grief, remorse, sadness
\item[anger] anger, annoyance, disapproval
\item[fear] fear, nervousness
\item[surprise] surprise, realization, confusion, curiosity
\item[disgust] disgust
\item[neutral] neutral
\end{description}
Family scores are the sums of the member probabilities, normalised to a distribution.

\section{Family centroids on the circumplex}
\begin{center}
\begin{tabular}{lrr}
\toprule
Family & Valence & Arousal \\
\midrule
joy & $+0.80$ & $0.60$ \\
sadness & $-0.75$ & $0.30$ \\
anger & $-0.70$ & $0.80$ \\
fear & $-0.65$ & $0.75$ \\
surprise & $+0.30$ & $0.80$ \\
disgust & $-0.70$ & $0.60$ \\
neutral & $0.00$ & $0.30$ \\
\bottomrule
\end{tabular}
\end{center}
Hand-set from the standard placement of the basic emotions on the circumplex; the minimum of Eq.~\ref{eq:va} over these points is $0.412$.

\section{Readable expression words}
A record is face-classified when the describer reports at least one person and its expression answer contains one of:
\begin{quote}
smile, smiling, happy, sad, angry, serious, surprised, surprise, calm, excited, crying, laughing, frowning, scared, shocked, worried, concentration, concentrating, grin, grinning, neutral, confused, determined, pensive, sleepy, sleeping.
\end{quote}
The people count is parsed from digits or number words (none, one, two, \dots, several, many, crowd, group, lots); a non-empty answer that is neither is read as one person present.

\section{Name redaction procedure}
\begin{enumerate}\setlength{\itemsep}{1pt}
\item \textbf{Detect.} Candidates from three detectors: a given-name gazetteer (7,944 names), mid-sentence capitalised tokens absent from a 236k-word English dictionary, and naming phrases (``called X'', ``my friend X''). High recall by design, because most prompts are written entirely in lower case.
\item \textbf{Review.} Every candidate has a recorded decision (redact or keep) with a reason and, for redactions, the record it applies to. The audit fails if any candidate is unreviewed or if a redaction entry matches nothing.
\item \textbf{Apply.} At load time, per record, replacing the token with \texttt{[name]}; consecutive placeholders collapse to one. The files on disk are never modified. Illness, grief and place names are kept.
\end{enumerate}

\section{Single-coder check: coding instructions}
The coder works from numbered contact sheets that show only the drawings, and fills four columns per drawing:
\begin{description}\setlength{\itemsep}{1pt}
\item[person] y / n --- is at least one person or face drawn?
\item[face] y / n --- is there a face whose expression can be read? (n if no person)
\item[valence] $+$ / $-$ / 0 --- the feeling the drawing gives: positive, negative, neither
\item[family] joy / sadness / anger / fear / surprise / disgust / neutral
\end{description}
The sample is fixed by seed (120 drawings: 40 face-classified, 40 scene-classified with no people reported, and every record with a vague or word-only people answer, topped up equally from the first two strata). The prompt, the description and the model readings are not shown during coding.
